
\documentclass{article}

% Sprach- und Zeichencodierung
\usepackage[utf8]{inputenc}
\usepackage[T1]{fontenc}
\usepackage[ngerman]{babel}

% Mathematische Symbole und Einheiten
\usepackage{amsmath}
\usepackage{amssymb}

% Seitenlayout
\usepackage[a4paper,margin=2.5cm]{geometry}

\begin{document}

\title{Analyse von Datenbank-Performance, Streuung und Indizierung}
\author{David Rinner}
\date{\today}

\maketitle

\section{Einleitung und Methodik}

Ziel dieser Analyse ist die Untersuchung des Verhaltens von
SQLite-Datenbanken bei großen Datenmengen ($\geq 500\,000$ Datensätze).
Dabei wurden die Streuung von Zufallsdaten, die mithilfe von Python und
\texttt{Faker} generiert wurden, die Auswirkungen von Indizes auf die
Performance sowie der Einfluss eines Daten-Bias untersucht.

\section{Streuungsanalyse der Zufallsdaten}

Da Namen keine Distanz im herkömmlichen Sinne besitzen, lässt sich ihre
Streuung nicht direkt über geometrische Abstände bestimmen.
Stattdessen wird die Streuung über die \textbf{Häufigkeitsverteilung der
Vornamen} definiert.

Jeder Vorname wird dabei als Kategorie betrachtet, deren Häufigkeit über
alle Datensätze hinweg gezählt wird.

\begin{itemize}
    \item \textbf{Durchschnittliche Häufigkeit:} $\approx 249{,}75$
    Vorkommen pro Vorname.
    
    \item \textbf{Varianz der Häufigkeiten:} $267{,}28$
    (Standardabweichung $\sigma \approx 16{,}35$).
\end{itemize}

Die geringe Varianz zeigt, dass die verwendete Bibliothek
\texttt{Faker} die Vornamen relativ gleichmäßig über die Datensätze
verteilt.

\section{Performance und Speicherbedarf mit und ohne Index}

Ohne Index muss die Datenbank bei einer Abfrage 
die gesamte Tabelle durchsuchen und benötigt in der vorliegenden
Messung etwa $0{,}178\,\mathrm{s}$.

Nach der Erstellung eines Indexes 
sinkt die Abfragezeit auf etwa $0{,}001\,\mathrm{s}$. Dadurch wird die
Abfrage in diesem Beispiel deutlich beschleunigt.

Dieser Geschwindigkeitsgewinn geht jedoch mit einem erhöhten
Speicherbedarf einher. Während die Datenbank vor der Erstellung des
Indexes $11\,755\,520\,\mathrm{Byte}$ 
beanspruchte, steigt der Speicherbedarf nach dem Aufbau des
B-Tree-Indexes auf $19\,578\,880\,\mathrm{Byte}$.

Der zusätzliche Speicherbedarf entsteht dadurch, dass neben der
eigentlichen Tabelle eine zusätzliche Indexstruktur gespeichert werden
muss.

\section{Untersuchung der Biased Database}

Wird die Verteilung der Vornamen gezielt verändert, sodass
50\,\% aller Datensätze denselben Vornamen (\texttt{Simon}) tragen,
verändert sich die Geschwindigkeit der Suche.

\begin{itemize}

    \item \textbf{Seltener Wert (\texttt{Peter}):}

    Wird nach einem selten vorkommenden Namen gesucht, kann der Index
    die passenden Datensätze sehr schnell finden. In der Messung wurde
    eine Abfragezeit von etwa $0{,}001\,\mathrm{s}$ erreicht.

    \item \textbf{Häufiger Wert (\texttt{Simon}):}

    Da \texttt{Simon} in 50\,\% der Datensätze vorkommt, muss die
    Datenbank sehr viele passende Datensätze finden. Dadurch ist die
    Suche mit dem Index langsamer als bei einem seltenen Namen.
    Mit Index benötigt die Abfrage etwa $0{,}050\,\mathrm{s}$,
    während ohne Index etwa $0{,}131\,\mathrm{s}$ benötigt werden.

\end{itemize}

\textbf{Grund für dieses Verhalten:}

Der Unterschied liegt darin, wie oft ein bestimmter Name in der
Datenbank vorkommt. Bei einem seltenen Namen wie \texttt{Peter} gibt
es nur wenige passende Datensätze. Der Index kann diese daher schnell
finden.

Bei \texttt{Simon} sieht es anders aus: Der Name kommt in der Hälfte
aller Datensätze vor. Bei insgesamt $500\,000$ Datensätzen sind das
ungefähr $250\,000$ Treffer. Der Index muss deshalb deutlich mehr
passende Einträge verarbeiten.

Trotzdem ist die Suche mit dem Index in diesem Beispiel noch schneller
als die Suche ohne Index. Das zeigt, dass ein Index besonders bei
selten vorkommenden Werten hilfreich ist. Bei sehr häufig vorkommenden
Werten fällt der Vorteil dagegen kleiner aus.


\section{Fazit}

Die Untersuchung zeigt, dass die Verwendung eines Indexes bei großen
Datenmengen einen erheblichen Einfluss auf die Abfragegeschwindigkeit
haben kann. Besonders bei Suchbegriffen 
mit selten vorkommenden Werten ist der Geschwindigkeitsgewinn deutlich.

Gleichzeitig erhöht ein Index den Speicherbedarf der Datenbank, da
zusätzliche Strukturen angelegt und verwaltet werden müssen.

Die Untersuchung der \texttt{bias.db} zeigt außerdem, dass die
Verteilung der Daten einen wesentlichen Einfluss auf die Effektivität
eines Indexes hat. Bei häufig vorkommenden Werten ist der Vorteil eines
Indexes geringer als bei seltenen Werten.

Damit wird deutlich, dass bei der Optimierung von Datenbanken nicht
nur die Anzahl der Datensätze, sondern auch deren Verteilung und die
Selektivität der verwendeten Suchbedingungen berücksichtigt werden
müssen.

\end{document}

